\chapter{Problem Definition}
\label{ch:problem}

\section{Introduction}
Chapter~\ref{ch:intro} introduced the project and Chapter~\ref{ch:literature}
surveyed the work it builds on. This chapter states the problem precisely. It
sets out what is wrong with the air-quality information currently available to
the Pakistani public, examines the systems that exist and where each of them
stops, distils the shortcomings into a numbered set of gaps, outlines the
solution this project proposes against each of them, records the assumptions and
constraints under which that solution was built, assesses its feasibility, and
enumerates the deliverables by which its completion can be judged.

\section{Problem Statement}
Air quality in Pakistan is a severe, recurring and largely predictable
public-health problem, and the information infrastructure available to respond to
it is inadequate in four distinct ways.

\textbf{It reports the present rather than the future.} Official and commercial
sources publish the Air Quality Index as it stands at a monitoring station. This
is a measurement, not a forecast, and it arrives too late to inform the decisions
people take: whether to walk a child to school tomorrow morning, whether to shift
an outdoor work detail to a cleaner day, whether an asthmatic should refill a
prescription before a smog episode. The World Health Organization's
guidelines~\cite{ref:who} are framed around exposure, and exposure is something
that can only be reduced by acting \emph{before} it occurs.

\textbf{It covers a few places and ignores the rest.} Where forecasts do exist,
they overwhelmingly concern Lahore and Karachi. The exploratory analysis
conducted for this project, presented in Section~\ref{sec:design-eda}, found that
Faisalabad's mean AQI over the study period was in fact \emph{higher} than
Lahore's, and that Sargodha, Multan, Gujranwala and Sialkot all sit above the
threshold the EPA labels unhealthy. Residents of these cities, and of the
smaller cities of Khyber Pakhtunkhwa, Balochistan, Gilgit-Baltistan and Azad
Jammu \& Kashmir, are served by nothing at all. A per-city modelling approach
scales badly to this problem, because each new city requires its own model, its
own training run and its own maintenance.

\textbf{It explains nothing and quantifies nothing.} A bare index value gives a
user no way to distinguish a genuine pollution episode from an unusual
meteorological configuration, no indication of how confident the source is, and
no means of judging whether the source has been right in the past. Trust in a
public-health signal is not established by assertion; it is established by
showing the working. The literature offers the tools - Shapley
attributions~\cite{ref:shap,ref:treeshap} for the first concern and conformal
intervals~\cite{ref:leiconformal} for the second - but they are seldom delivered
to an end user.

\textbf{It is maintained by hand, if at all.} This is the deepest of the four
problems, and the one the machine-learning literature most often overlooks. Air
quality is not stationary: the relationship between weather, calendar and
pollution shifts with the season, with agricultural cycles and with changes in
emission. A model trained once and left alone therefore decays, and the decay is
silent - the model continues to emit confident numbers that are progressively
less correct. Sculley et al.~\cite{ref:techdebt} identified precisely this class
of failure as the dominant source of technical debt in deployed machine-learning
systems, and Gama et al.~\cite{ref:driftsurvey} characterised the drift mechanism
behind it. A forecasting service for a seasonal phenomenon must therefore be
self-maintaining, or it will quietly stop working within a year of deployment.

Stated formally: \emph{there is a need for an automated system that forecasts the
Air Quality Index several days ahead for many Pakistani cities rather than one,
that retrains and validates itself without human intervention while refusing to
deploy a model that is not measurably better than the one in service, that
quantifies its own uncertainty and monitors its own inputs and outputs for decay,
that explains each individual prediction in language a non-specialist can act on,
and that is delivered publicly and reliably at negligible cost.}

\section{Existing Systems and Current Solutions}
Four families of existing system bear on this problem. Each solves part of it and
stops somewhere specific.

\subsection{Government and Regulatory Monitoring}
Provincial environmental protection agencies operate reference-grade monitoring
stations and publish index readings. These are the authoritative source for
\emph{what the air is now}, and nothing in this project attempts to displace them.
Their limitations for the present purpose are that the station network is sparse
relative to the population it serves, that the published output is a nowcast, and
that no explanation or forecast accompanies it.

\subsection{Consumer Air-Quality Applications}
A number of commercial applications and websites aggregate station and satellite
data into a consumer-facing index, and the better ones include a short forecast.
Their weaknesses, from the standpoint of this problem, are that the model behind
the forecast is undisclosed and unevaluated - no per-horizon accuracy is
published, so a user cannot know whether the third day of a three-day forecast
carries any skill at all - that no explanation is offered, and that coverage of
smaller Pakistani cities is thin. They are products rather than
scientific instruments, and they are not designed to be checked.

\subsection{Physical Chemical Transport Models}
Services such as the Copernicus Atmosphere Monitoring System~\cite{ref:cams}
produce genuinely predictive composition forecasts from first principles. They
are the correct tool for scientific attribution and for regional policy analysis.
They are also global-scale scientific infrastructure: they presuppose an
emissions inventory, meteorological boundary conditions and substantial computing
capacity. They cannot be reproduced by an individual developer, they are not
tuned to individual Pakistani cities, and their output is not delivered in a form
a resident can act on.

\subsection{Published Machine-Learning Studies}
The academic literature on machine-learning air-quality forecasting is
large~\cite{ref:raq,ref:xgbpm,ref:lstmpm}, and it establishes convincingly that
supervised models can forecast pollutant concentrations with useful skill. Its
characteristic limitation is that it addresses the modelling question and stops.
A typical paper reports a single averaged accuracy figure for a model trained
once on a fixed dataset, frequently for one city, sometimes on a randomly
partitioned validation set that allows the model to train on the future. There is
usually no retraining pipeline, no promotion gate, no durable registry, no drift
instrumentation, no deployed interface and no ongoing measurement of whether the
model is still right. The model is a result; it is not a service.

\section{Identified Problems and Gaps}
The preceding survey reduces to seven concrete gaps, each of which is addressed
by a specific component of the system described in this thesis.

\begin{xltabular}{\textwidth}{Q{1.3cm} P{4.6cm} Y}
\caption{Identified gaps and the component of this system that addresses each.}\label{tab:gaps}\\
\toprule
\bfseries ID & \bfseries Gap in current practice & \bfseries How this project addresses it \\
\midrule
\endfirsthead
\toprule
\bfseries ID & \bfseries Gap in current practice & \bfseries How this project addresses it \\
\midrule
\endhead
\bottomrule
\endlastfoot
G1 & Information is a nowcast, so it arrives after the decision it should have
informed. & A 72-hour hourly forecast, refreshed daily, aggregated into both an
hourly curve and a three-day view. \\
G2 & Coverage is limited to one or two major cities; per-city models scale
badly. & A single global model conditioned on location, trained across 22 cities
covering every province, the capital territory, Gilgit-Baltistan and Azad Jammu
\& Kashmir. \\
G3 & Accuracy is unreported, or reported as one averaged number that conceals
decay with lead time. & Per-horizon RMSE, MAE and $R^2$ against a per-horizon
persistence baseline, plus a five-fold walk-forward backtest, published in the
interface. \\
G4 & Point forecasts are presented without uncertainty. & An 80\% prediction
interval derived from empirical residual quantiles, widening with horizon, drawn
as a band on every chart. \\
G5 & Predictions are unexplained, so they cannot be interrogated or trusted. &
Per-forecast SHAP attributions rendered as a plain-language sentence, plus an
interactive what-if simulator over the model's own drivers. \\
G6 & Models are trained once and decay silently as the relationship shifts. &
An hourly feature pipeline and a nightly retraining pipeline, with a
champion-challenger gate that promotes only a measurable improvement and records
every decision. \\
G7 & Nothing watches the system after deployment. & Population Stability Index
drift monitoring per feature, plus a forecast log joined back to realised
outcomes so that live error is measured rather than assumed. \\
\end{xltabular}

\section{Proposed Solution Overview}
The proposed solution is an operating system rather than a model, built on the
Feature-Training-Inference pattern~\cite{ref:fti} and running entirely on free
scheduled infrastructure. Its shape is as follows.

An \textbf{hourly feature pipeline} fetches the last seven days of pollutant and
weather observations for all 22 cities from a free, keyless public
API~\cite{ref:openmeteo}, engineers a 74-column feature table for each city
independently, and upserts the result into a managed feature store keyed on city
and timestamp. A one-off \textbf{backfill} runs the same logic over the full
history from January 2023 to the present, in three-month chunks, so that the
store contains roughly 697,000 hourly rows before the first model is trained.

A \textbf{daily training pipeline} reads the entire feature table, builds a
multi-horizon supervised dataset in which the forecast horizon is itself an input
feature, splits it chronologically, trains four candidate models - a persistence
baseline, a regularised linear model, a random forest and a gradient-boosted
ensemble - and scores each on a validation window that lies strictly in the
future relative to the training window. The best fitted candidate becomes the
\emph{challenger}, and is promoted to \emph{champion} only if it improves on the
incumbent's validation RMSE. Every run, promoted or refused, is recorded in a
persistent leaderboard, and every promoted champion is written to a durable,
versioned model registry.

A \textbf{hourly batch inference pipeline} loads the reigning champion from that
registry, fetches each city's current pollution state and the coming four days of
forecast weather, assembles one model input per future hour, predicts the whole
72-hour curve, widens each point by the residual quantiles recorded for its
horizon, explains the peak hour with SHAP, raises a hazard alert if the peak
crosses the unhealthy threshold, and writes a single forecast document.

An \textbf{observability layer} computes the Population Stability Index for each
monitored feature between a recent window and the historical reference, appends
every issued forecast to a log, and joins that log back to the air quality that
subsequently arrived so that realised error can be measured directly.

A \textbf{delivery layer} serves the forecast document, the leaderboard, the
evaluation results and the monitoring report through a typed HTTP API, presents
them in a five-tab browser dashboard with an interactive map of Pakistan, and
exposes the same grounded data to a language-model advisor through a standard
tool protocol~\cite{ref:mcp}.

\section{Assumptions and Constraints}

\subsection{Assumptions}
\begin{enumerate}
  \item \textbf{The data source is representative and available.} Open-Meteo's
  air-quality product is derived from the Copernicus atmospheric composition
  reanalysis and forecast, and its weather archive from ERA5~\cite{ref:era5}.
  These are modelled rather than station-measured values for a given city, and
  the system assumes they are an adequate proxy for the air a resident of that
  city breathes.
  \item \textbf{The weather forecast is usefully accurate.} At training time the
  model sees the weather that actually occurred at the target hour; at inference
  time it sees the forecast weather for that hour. The system assumes that the
  gap between the two is small enough not to invalidate the learned relationship.
  This is standard practice in operational forecasting, but it means the reported
  validation error is a slight under-estimate of true operational error, and
  Chapter~\ref{ch:testing} treats that honestly.
  \item \textbf{The target definition is stable.} The EPA index definition and
  its category boundaries~\cite{ref:epa} are assumed not to change during the
  system's life.
  \item \textbf{Recent history is available at inference time.} The anchor state
  - the observed pollution at the moment the forecast is issued - is assumed to
  be retrievable for every city at every run.
\end{enumerate}

\subsection{Constraints}
\begin{enumerate}
  \item \textbf{Zero budget.} Every component must sit within a free tier. This
  ruled out managed training infrastructure, paid data, always-on servers and
  GPU compute, and it directly shaped several design decisions documented in
  Chapter~\ref{ch:impl}, including the decision to disable feature-store
  statistics after the profiler exhausted the free executor's memory.
  \item \textbf{Ephemeral compute.} Scheduled runners are destroyed when a job
  ends, so no artefact may be left on local disk if it is needed later. This is
  the constraint that forces a genuine model registry rather than a saved file.
  \item \textbf{Job time limits.} A scheduled job must finish inside the
  platform's cap, which constrains the size of the training set and forced the
  backfill to be chunked, resumable and non-blocking.
  \item \textbf{Development platform.} Development was carried out on Windows,
  where the feature-store client cannot be installed because one of its
  transitive cryptographic dependencies fails to build. All feature-store work
  therefore runs on Linux continuous integration, and local development uses an
  interchangeable Parquet backend behind the same interface.
  \item \textbf{Source granularity.} The finest resolution available is one city
  and one hour; street-level or sub-hourly forecasting is therefore impossible
  regardless of modelling choices.
  \item \textbf{Public exposure.} The deployed interface is public and
  unauthenticated, which constrains it to read-only, non-sensitive data and makes
  cost predictability a design concern.
\end{enumerate}

\section{Feasibility Analysis}

\subsection{Technical Feasibility}
The technical risks were identified at the outset and each was resolved by a
concrete mechanism.

The first risk was whether a free-tier feature store could absorb a
three-and-a-half-year, 22-city backfill. It could, but not naively: a blocking
insert waited indefinitely on a stalled materialisation job and froze the run at
the tenth city. Converting the write to a non-blocking submission, writing city by
city rather than in one batch, and making the resume logic history-aware rather
than presence-aware resolved it, and the resulting design is idempotent, so a
failed run can simply be repeated.

The second risk was whether a single global model could serve 22 cities at ten
horizons without being worse than a per-city model at any of them. The evidence
in Chapter~\ref{ch:testing} is that it can: the global model beats the
persistence baseline at every horizon from two hours to 72, by margins of between
16\% and 41\%, and it retains a substantial margin at the long horizons where a
local or univariate approach is weakest. The one exception, a one-hour lead time
at which the trivial predictor is better, is reported and explained rather than
omitted.

The third risk was whether the daily retraining job would fit inside the
scheduler's time limit. Subsampling the supervised set to 200,000 rows brings a
full four-candidate training run comfortably inside the cap, and this is the
principal reason the gradient-boosted model rather than the recurrent model is
the production choice: it trains in about a minute against roughly twenty-five
for the sequence model on the same hardware.

The fourth risk was whether the trained model could be served without a live
feature-store connection, since the deployment image deliberately omits that
dependency. It can, because the inference stage writes a self-contained forecast
document and the model bundle is shipped with the image.

\subsection{Economic Feasibility}
The system's recurring cost is essentially zero, which was a requirement rather
than an accident. Data acquisition is free and requires no key. Scheduled compute
uses the free allowance of the repository host's continuous-integration service;
the hourly feature job and the nightly training job together consume a small
fraction of it. Feature storage and the model registry use a free managed tier.
The frontend is served from a free static edge host. The backend runs as a small
container on a low-cost container host. The only component with a genuinely
variable cost is the optional language-model advisor, which is metered per
request and degrades gracefully - returning a clear message rather than an error
- when no key is configured. The consequence is that the service can keep
running indefinitely without a funding decision, which is precisely what a
public-interest tool needs.

\subsection{Operational Feasibility}
Operationally, the system is designed to require no one. Ingestion, retraining,
evaluation, inference and monitoring are all scheduled; the refreshed artefacts
are committed automatically; the frontend reads them without redeployment. The
promotion gate means that an unattended retrain cannot degrade the service, since
a challenger that fails to improve is simply refused and the incumbent continues.
Each pipeline logs its decisions in a form an operator can audit after the fact,
and the leaderboard, evaluation report and monitoring report are all published in
the interface rather than buried in logs. The realistic operational burden is
periodic review rather than daily intervention.

\section{Deliverables and Development Requirements}
Table~\ref{tab:deliverables} lists the deliverables against which the project's
completion is assessed.

\begin{xltabular}{\textwidth}{Q{1.3cm} P{4.9cm} Y}
\caption{Project deliverables.}\label{tab:deliverables}\\
\toprule
\bfseries ID & \bfseries Deliverable & \bfseries Description \\
\midrule
\endfirsthead
\toprule
\bfseries ID & \bfseries Deliverable & \bfseries Description \\
\midrule
\endhead
\bottomrule
\endlastfoot
D1 & Ingestion and feature engineering & A retrying client for three Open-Meteo
endpoints, an EPA index implementation, and a per-city feature builder producing
74 columns from 21 raw ones. \\
D2 & Populated feature store & Roughly 697,000 hourly rows for 22 cities from
January 2023 onward, upserted on a composite key, with a resumable backfill. \\
D3 & Exploratory analysis & City ranking, seasonality, diurnal profile, weather
correlation and distribution figures over the full record. \\
D4 & Multi-horizon supervised builder & Assembly of target-time, anchor-state and
horizon features with a chronological split and unit tests asserting the absence
of leakage. \\
D5 & Training pipeline & Four candidate families, honest validation, residual
quantile intervals, and a bundled artefact carrying the estimator, column order,
interval table and metrics. \\
D6 & Champion-challenger gate and registry & A persistent leaderboard, a
promotion rule, and a durable versioned model registry from which inference
loads. \\
D7 & Evaluation suite & Per-horizon metrics against a per-horizon baseline, a
five-fold walk-forward backtest, and a published report and figures. \\
D8 & Sequence-model study & An LSTM trained and evaluated on the same task and
split, with a reasoned promotion decision. \\
D9 & Inference pipeline & A 72-hour forecast with intervals, categories, alerts
and explanations for every city, written as a single document. \\
D10 & Explainability & TreeSHAP attributions rendered in plain language,
per-forecast in the interface and globally as an importance figure. \\
D11 & What-if simulator & Live re-prediction under user-overridden drivers,
presented as baseline versus scenario. \\
D12 & Monitoring & Per-feature PSI drift with status thresholds, a forecast log,
and realised error scored against observed outcomes. \\
D13 & API and dashboard & A typed HTTP interface and a five-tab React
application with charting, cartography and alerting. \\
D14 & Advisor and tool server & Four grounded tools exposed both to an in-app
advisor and over the Model Context Protocol. \\
D15 & Continuous integration and deployment & Three scheduled workflows, a
container image for the backend, a static build for the frontend, and a live
public deployment. \\
D16 & Tests and documentation & Unit tests over the critical logic, and this
thesis. \\
\end{xltabular}

\section{Summary of Problem Definition}
This chapter defined the problem the \projname\ addresses: air-quality
information for Pakistan that is current-only, geographically narrow,
unexplained, unquantified and manually maintained. It examined the four families
of existing system - regulatory monitoring, consumer applications, physical
transport models and published machine-learning studies - and identified where
each stops. It distilled seven gaps, outlined the pipeline-based solution
proposed against them, stated the assumptions and the free-tier constraints under
which that solution was built, argued its technical, economic and operational
feasibility, and enumerated sixteen deliverables. The next chapter converts this
into a formal requirement specification.
